\documentclass[xcolor=dvipsnames,aspectratio=169]{beamer}

\usepackage{fontspec}
\defaultfontfeatures{Ligatures=TeX}
%\setsansfont{Liberation Sans}
\usepackage{polyglossia}
\setdefaultlanguage{german}

\usetheme{Berlin}
\usecolortheme[named=LimeGreen]{structure}
\usepackage{beamerthemesplit} % kam neu dazu
	
\usepackage{amsmath}
\usepackage{graphicx}
\graphicspath{{pictures/}}
\usepackage{amssymb}
\usepackage{amsfonts}
\usepackage{caption}
\usepackage{multimedia}
\usepackage{tikz}
\usepackage{listings}
\usepackage{acronym}
\usepackage{uhrzeit}
\usepackage{lmodern}
\usepackage{multicol}


\definecolor{pblue}{rgb}{0.13,0.13,1}
\definecolor{pgreen}{rgb}{0,0.5,0}
\definecolor{pred}{rgb}{0.9,0,0}
\definecolor{pgrey}{rgb}{0.46,0.45,0.48}

\newcounter{countitems}
\newcounter{nextitemizecount}
\newcommand{\setupcountitems}{%
  \stepcounter{nextitemizecount}%
  \setcounter{countitems}{0}%
  \preto\item{\stepcounter{countitems}}%
}
\makeatletter
\newcommand{\computecountitems}{%
  \edef\@currentlabel{\number\c@countitems}%
  \label{countitems@\number\numexpr\value{nextitemizecount}-1\relax}%
}
\newcommand{\nextitemizecount}{%
  \getrefnumber{countitems@\number\c@nextitemizecount}%
}
\newcommand{\previtemizecount}{%
  \getrefnumber{countitems@\number\numexpr\value{nextitemizecount}-1\relax}%
}

\makeatother    
\newenvironment{AutoMultiColItemize}{%
\ifnumcomp{\nextitemizecount}{>}{3}{\begin{multicols}{2}}{}%
\setupcountitems\begin{itemize}}%
{\end{itemize}%
\unskip\computecountitems\ifnumcomp{\previtemizecount}{>}{3}{\end{multicols}}{}}

\lstset{
    escapeinside={(*}{*)}
}

\lstdefinestyle{Java}{
  showspaces=false,
  showtabs=false,
  tabsize=2,
  breaklines=true,
  showstringspaces=false,
  breakatwhitespace=true,
  commentstyle=\color{pgreen},
  keywordstyle=\color{pblue},
  stringstyle=\color{pred},
  basicstyle=\footnotesize\ttfamily,
  numbers=left,
  numberstyle=\tiny\color{gray}\ttfamily,
  numbersep=7pt,
  %moredelim=[il][\textcolor{pgrey}]{$$},
  moredelim=[is][\textcolor{pgrey}]{\%\%}{\%\%},
  captionpos=b
}

\lstdefinestyle{basic}{  
  basicstyle=\footnotesize\ttfamily,
  breaklines=true
  numbers=left,
  numberstyle=\tiny\color{gray}\ttfamily,
  numbersep=7pt,
  backgroundcolor=\color{white},
  showspaces=false,
  showstringspaces=false,
  showtabs=false,
  frame=single,
  rulecolor=\color{black},
  captionpos=b,
  keywordstyle=\color{blue}\bf,
  commentstyle=\color{gray},
  stringstyle=\color{green},
  keywordstyle={[2]\color{red}\bf},
}


\lstdefinelanguage{custom}
{
morekeywords={public, void},
sensitive=false,
morecomment=[l]{//},
morecomment=[s]{/*}{*/},
morestring=[b]",
}


\lstdefinestyle{BashInputStyle}{
  language=bash,
  showstringspaces=false,
  basicstyle=\small\sffamily,
  numbers=left,
  numberstyle=\tiny,
  numbersep=5pt,
  frame=trlb,
  columns=fullflexible,
  backgroundcolor=\color{gray!20},
  linewidth=0.9\linewidth,
  xleftmargin=0.1\linewidth
}

%Logo in the upper right just change if you know what you are doing^^
\addtobeamertemplate{frametitle}{}{%
\begin{tikzpicture}[remember picture,overlay]
\node[anchor=north east,yshift=2pt] at (current page.north east) {\includegraphics[height=1.8cm]{htw}};
\end{tikzpicture}}

\begin{document}
\bibliographystyle{alpha}
\title{Netzwerke -- Übung Wintersemester 2020/21}
\subtitle{Organisatorisches\\

\href{mailto:Benjamin.Troester@HTW-Berlin.de}{Benjamin.Troester@HTW-Berlin.de}\\
		PGP: ADE1 3997 3D5D B25D 3F8F 0A51 A03A 3A24 978D D673 }

\author{Benjamin Tröster}

\date{}

\begin{frame}
\titlepage
\end{frame}

\section*{Road-Map}
\begin{frame}
\frametitle{Road-Map}
\begin{multicols}{2}
  \tableofcontents
\end{multicols}
\end{frame}

\section{Disclamer}
\begin{frame}
	\frametitle{Disclamer I}
	\begin{itemize}
		\item Objektive Kritik ist immer willkommen!
		\item Wenn sie Fehler oder Anmerkungen haben, tragen sie diese ruhig vor.
		\item Sie wollen ehrliches Feedback genauso wie ich
		\item Wenn Sie Anregungen, Verbesserungsvorschläge haben, nehme ich diese ebenfalls gerne an.
	\end{itemize}
\end{frame}

\section{Organisatorisches}
\subsection{Übungsablauf}
\begin{frame}
	\frametitle{Übungsablauf}
	\begin{itemize}
		\item Voraussichtlich ab dem 02.11.2020 via Distanzlehre
		\item Meeting/Übung via Moodle: Telco-Plugin \emph{BigBlueButton (BBB)}
		\item Übungsgruppe \& -zeit $\Rightarrow$ LSF
		\item Übungsblätter, Folien, Literatur \& Links etc. $\Rightarrow$ \url{moodle.htw-berlin.de}
		\item Hausaufgaben $\rightarrow$ Vorbereitung der kommenden Übung
		\item Laboraufgaben $\rightarrow$ Anwendung der Hausaufgaben
		\item Zweiwöchentlicher Tonus 
		\begin{itemize}
			\item \textasciitilde 5-6 Übungsblätter -- bestehend aus Theorie- \& Laborteil
			\item Jede Übung in sich abgeschlossen
			\item \textbf{Aber:} Übungen bauen aufeinander auf
		\end{itemize}
	\end{itemize}
\end{frame}


\subsection{ECTS}
\begin{frame}
	\frametitle{Arbeitsaufwand -- ECTS}
	\begin{itemize}
		\item Modul Netzwerke: 5 ECTS (European Credit Transfer System) -- manchmal auch LP/CP (Leistungspunkte)
			\begin{itemize}
				\item 1 ECTS $\widehat{=}$ 30h
				\item Workload Netzwerke: 150h/Semester oder 37,5h/Monat oder $\sim$ 9,375h/Woche
				\item 2 SWS \footnote{Semesterwochenstunden} VL + 2 SWS Übung $\rightarrow$ 4 SWS oder 3h
				\item D.h. restlichen \textbf{6h/Woche}
				\begin{itemize}
					\item Vorbereitung \& Nachbereitung der Vorlesung
					\item Hausaufgaben, Recherche...
				\end{itemize}
			\end{itemize}
			\item Total Workload pro Semester: 30 CP $\widehat{=}$ 900h $\rightarrow$  $\sim$56.25/Woche
	\end{itemize}
\end{frame}

\begin{frame}
	\frametitle{Selbstverantwortliches Lernen \textbf{Wichtig}}
	\begin{itemize}
		\item Sie sind für den Fortschritt ihrer Lernziele selbst verantwortlich
		\item Im eigenen Interesse:
		\begin{itemize}
			\item Bearbeiten Sie die Hausaufgaben
			\item Notieren Sie sich fragen
			\item Suchen Sie nach Material, das Sie beim lernen unterstützt
		\end{itemize}
		\item Unterstützungsangebot:
		\begin{itemize}
			\item Literatur \& Links auf Moodle
			\item Bilden Sie Übungsgruppen
			\item Plenum -- basierend auf Ihrer Vorbereitung
			\begin{itemize}
				\item Fragen bitte im Moodle-Forum vorher stellen!
			\end{itemize}
		\end{itemize}
		\item Meine Aufgabe ist es Sie bei der Umsetzung der Übungen und dem Verstehen des Kursinhaltes zu unterstützen!
		\begin{itemize}
			\item Sie sollen eigene Lösungswege finden!
			\item Fehler machen und daraus lernen!
		\end{itemize}
	\end{itemize}
\end{frame}


\section{Motivation}
\subsection{Studien- \& Prüfungsordnung AMB30/2012}
\begin{frame}
	\frametitle{Motivation}
	Studien- \& Prüfungsordnung AMB30/2012 S. 375\\
	B12 Netzwerke
	\begin{itemize}
		\item Die Studierenden erwerben Kenntnisse wichtiger Netzwerkprotokolle und -dienste
		\item Sie erwerben Fertigkeiten im Aufbau von Rechnernetzen
		\item Sie erwerben Grundlagen sicherer Netzwerkkommunikation
		\item Sie können einfache Netzwerke realisieren
	\end{itemize}
\end{frame}
\begin{frame}
	\frametitle{Motivation}
	\vspace{-0.7cm}
	Etwas genauer:\\
	Studienordnung des Bachelorstudiengangs Angewandte Informatik -- B12 Netzwerke:
	\begin{itemize}
		\item Die Studierenden erwerben Kenntnisse wichtiger Netzwerkprotokolle und -dienste
		\item Sie erwerben Fertigkeiten im Aufbau von Rechnernetzen
		\item Sie erwerben Grundlagen sicherer Netzwerkkommunikation
		\item Sie können einfache Netzwerke realisieren
		\item Grundlagen der Systemverwaltung
		\item OSI-Referenzmodell
		\item Netzwerkprotokolle TCP, UDP, IP, ...
		\item Routing
		\item Name Service
		\item HTTP
	\end{itemize}
\end{frame}

\subsection{Netzwerktechnik in anderen Modulen}
\begin{frame}
\frametitle{Motivation II}
	\begin{itemize}
		\item Netzwerke $\rightarrow$ Fundament der digitalen Infrastruktur
		\item Grundlagen der Veranstaltungen Netzwerke sind Grundlagen anderer Kurse
		 \begin{AutoMultiColItemize}
 			\item[1.] Netzwerke
 			\item[2.] Betriebssysteme -- Client-Server-Programmierung, Sockets, Grid-Computing...
 			\item[3.] Programmieren 3 \& Datenbanken -- Sockets, JDBC...
 			\item[4.] Webentwicklung \& Verteilte Systeme -- MVC-Pattern, Messaging...
 			\item[5.] KBE, Spez. Anwd. \& Projektstudium -- JSF, JPA, CDI...
 			\item[6.] Möglicherweise Abschlussarbeit
  		\end{AutoMultiColItemize}
	\end{itemize}
\end{frame}

\end{document}